\section{本論文の構成と前提知識}
\label{sec:organization-prerequisites}

本論文では，複素ベクトル空間，内積，正規直交基底，
線形写像，固有値・固有ベクトル，ユニタリ行列，エルミート行列，
およびスペクトル分解に関する線形代数の基本事項を既知とする．
これらを学んだ学生を読者として想定し，量子計算の知識は仮定しない．
一方，量子計算で用いるブラケット記法，量子状態と測定の公理，
量子回路の表記については本文中で定義する．
また，テンソル積は複数量子ビットの記述とHHLアルゴリズムの
レジスタ構成に不可欠であるため，その定義と必要な性質を第2章で詳しく述べる．

本論文を通して，対象とする量子系は有限次元とし，
特に断らない限り純粋状態を扱う．
ベクトルを量子状態として表すときにはケット$\ket{\psi}$を用い，
その共役転置をブラ$\bra{\psi}$と書く．
内積$\braket{\phi|\psi}$は第1変数について共役線形，
第2変数について線形とする．
内積から定まるノルムを$\lVert\cdot\rVert$，
行列または線形写像の随伴を$A^\dagger$によって表す．
これらの記法の具体的な性質は，第2章で量子状態を導入する際に確認する．

各章の内容と，HHLアルゴリズムにおける役割を
表\ref{tab:introduction-organization}に示す．
\begin{table}[H]
  \centering
  \caption{本論文の構成とHHLアルゴリズムとの関係}
  \label{tab:introduction-organization}
  \begin{tabular}{c|>{\raggedright\arraybackslash}p{4.2cm}|>{\raggedright\arraybackslash}p{7.6cm}}
    \hline
    章 & 主な内容 & HHLアルゴリズムにおける役割 \\
    \hline
    第2章
      & 量子状態，テンソル積，量子ゲート，測定，補助量子ビット，データの符号化，制御回転
      & 入力状態と複数レジスタを表し，固有値に応じた振幅の変換と測定後状態を計算するための基礎を与える． \\
    第3章
      & 量子フーリエ変換，量子位相推定，振幅増幅，行列の指数関数
      & 行列の固有値を位相として取り出し，HHLの成功確率を改善するためのサブルーチンを与える． \\
    第4章
      & 量子線形方程式問題，HHLの構成と正しさ，具体例，成功確率，有限精度と成立条件
      & 第2章と第3章の結果を組み合わせ，成功時に解状態が得られることを証明する． \\
    第5章
      & 結果と成立条件の整理
      & 本論文で示した結果と，計算量を評価するために必要な仮定をまとめる． \\
    \hline
  \end{tabular}
\end{table}

第2章では，量子計算の基礎を扱う．
量子ビットと量子状態から始め，量子状態の時間発展を表すユニタリ変換，
テンソル積による複合系，量子ゲートと量子回路，射影測定を順に導入する．
さらに，補助量子ビットと逆演算，古典ベクトルの振幅符号化，
レジスタの値に応じた制御回転を説明する．

第3章では，HHLアルゴリズムに必要な量子サブルーチンを扱う．
まず量子フーリエ変換を定義し，ユニタリなゲートによる回路分解を示す．
次に，量子位相推定がユニタリ行列の固有値の位相を
量子レジスタへ符号化する過程を導出する．
その後，成功状態の振幅を増やす振幅増幅を二次元空間上の回転として説明する．
最後に，エルミート行列$A$からユニタリ行列$e^{\mathrm{i}At}$を構成し，
$A$の固有値と位相推定で得られる位相の関係を示す．

第4章では，量子線形方程式問題を定義し，HHLアルゴリズムを構成する．
位相推定，制御回転，逆位相推定，測定の各段階について，
三つの量子レジスタの状態を順に計算する．
その上で，理想的な場合の出力状態が
式\eqref{eq:introduction-solution-state}に一致することを証明する．
また，$2\times2$行列の例を用いて全過程を具体的に計算し，
成功確率と振幅増幅，有限精度による誤差，計算量を述べる際の仮定を整理する．

第5章では，各章で得られた結果をまとめ，
HHLアルゴリズムの正しさが成立する数学的仮定と，
計算上の利点を論じるために別途必要となる実装上の仮定を区別して述べる．

本論文の範囲は，HHLアルゴリズムの理解に必要な事項に限る．
量子計算の歴史，量子ハードウェア，ブロッホ球，量子テレポーテーション，
密度行列，振幅推定，変分量子アルゴリズム，
およびHHL以後の個別の量子機械学習アルゴリズムは扱わない．
